\documentclass[12pt]{article}
\usepackage[utf8]{inputenc}
\usepackage[T1]{fontenc}
\usepackage{hyperref}
\title{Northwestern Fars Language Resources for K-12 Curriculum Development}
\author{Compiled by Claude Sonnet 5.5 (claude-sonnet-5-5) for Andr\'as Kornai}
\date{}
\begin{document}
\maketitle

\section{Introduction}
The code \texttt{faz} (Northwestern Fars) was assigned to a supposed language of Iran. The Ethnologue entry states that its editors could not identify any evidence that the language exists and were moving to have the code removed from ISO 639. I did not confirm whether the removal has been completed. No educational resources can be cited.

\section{Primary Educational Resources}
None found.

\section{Linguistic and Reference Materials}
\begin{itemize}
\item Ethnologue entry for faz: \url{https://www.ethnologue.com/language/faz}
\item Glottolog entry nort2640 (family ``Bookkeeping''): \url{https://glottolog.org/resource/languoid/id/nort2640.bigmap.html}
\end{itemize}

\section{Religious and Community Resources}
None found.

\section{Digital and Audio Resources}
None found.

\section{Curriculum Development Context}
There is no evidence of a speech community, so no curriculum can be developed for this code. Speakers of Fars-area varieties should be looked for under other codes.
\end{document}